Robotic manipulation of deformable materials is challenging because the physical properties governing material response---stiffness, cohesion, plasticity---are hidden and vary substantially across material families. We propose \PTA{}, a deliberately minimal two-stage framework in which the robot executes a short, fixed probing sequence, encodes the resulting proprioceptive and aggregate-particle trace into a latent task-conditioning vector, and conditions a residual manipulation policy on that vector. We evaluate \PTA{} on a reproducible cross-material edge-push benchmark in the Genesis multi-physics simulator spanning sand, snow, and elastoplastic media; all methods are trained on sand only and evaluated zero-shot on the others. The empirical picture is asymmetric and conditional: on the out-of-distribution elastoplastic split, \PTA{} improves transfer efficiency by 14.7 percentage points (60.7\% vs.\ 46.0\%) and reduces spillage by 14.7 percentage points (39.3\% vs.\ 54.0\%) over a matched reactive PPO baseline, with matched ablations indicating that both the fixed probe and the latent conditioning are necessary for this elastoplastic improvement; on the four other splits (in-distribution sand, two parameter-shifted sand variants, cohesive snow), \PTA{} underperforms the reactive baseline by 13--22 percentage points. We do not claim broad cross-material robustness. We offer the asymmetry as evidence consistent with a \emph{recoverable-deformation hypothesis}: probing appears beneficial when probe-induced perturbations relax on the task timescale and detrimental when they do not. A complementary observation---a privileged-parameter baseline that receives ground-truth material parameters collapses to 0\% transfer on elastoplastic across all three seeds---suggests that, in this sand-only-trained matched setup, an in-episode interaction trace was a more useful conditioning signal than a passive parameter vector. \PTA{} is therefore best read as a conditional tool for recoverable-response settings, not as a universal robustness method.
